% ============================================================================
%  Capitolo 12 -- Verm1, parte 2: la guida, il cibo e un punteggio (Lezione 12)
% ============================================================================
\chapter{Verm1, parte 2: la guida, il cibo e un punteggio}
\label{ch:verm1p2}

\epigraphline{Un gioco è una serie di scelte interessanti.}{Sid Meier}

\begin{chapgoals}
Alla fine di questo capitolo avrete
\begin{itemize}
\item la testa di un verme che guidate con il joystick, in quattro direzioni;
\item una regola che vieta l'inversione a U, scritta come uno XOR e un confronto;
\item il movimento preso da due piccole tabelle invece che da quattro salti;
\item il cibo che compare in una cella casuale, viene mangiato, e ricompare;
\item un punteggio, mostrato come tacche nella riga del titolo, e un suono quando il verme mangia.
\end{itemize}
\end{chapgoals}
% Generated by tools/textbook/export_listings.js from tools/decks/mastery.js and tools/concepts/registry.json: do not edit.
\begin{chapstatus}{Stato di apprendimento di questo capitolo}
\masterychip{now}{DA COMPRENDERE OGGI}\enspace le maschere di bit; i codici; gli array in memoria; gli ingressi con pull-up; l'algebra di Boole; la virgola fissa; i multiplexer; istruzioni con o senza segno; comparatori e parità; gli argomenti in a0 ... a7.\par\smallskip
\masterychip{use}{DA UTILIZZARE OGGI}\enspace la convenzione di chiamata (Capitolo~\ref{ch:processor}) e, usati senza essere studiati nel corso, gli strumenti di PRG32, le chiamate del runtime PRG32, il debugger di PRG32-QT.\par\smallskip
\masterychip{later}{DA COMPRENDERE PIÙ AVANTI}\enspace niente.
\end{chapstatus}
\begin{chapstatus}{Percorso del capitolo}
Alla fine del capitolo saprete:
\begin{enumerate}[nosep,leftmargin=2em]
\item leggere i pulsanti con le maschere e rifiutare un'inversione a U con un solo XOR
\item confrontare posizioni con XOR e OR per riconoscere il cibo
\item tenere e mostrare un punteggio
\end{enumerate}
Si appoggia su: le maschere di bit (Capitolo~\ref{ch:boole1}); i codici (Capitolo~\ref{ch:combinational}); gli array in memoria (Capitolo~\ref{ch:information}); gli ingressi con pull-up (Capitolo~\ref{ch:electronics}); l'algebra di Boole (Capitolo~\ref{ch:boole1}). Prossimo passo: Capitolo~\ref{ch:sequential}.
\par\smallskip\textbf{Guardando indietro, guardando avanti.} Nella Lezione 07 una funzione che ne chiama un'altra salvava ra in un frame di 16 byte. Oggi le funzioni di Verm1 si chiamano più in profondità: ciascuna salva ra una volta, all'ingresso, e lo ripristina prima di ret. La Lezione 18 enuncia le regole: quali registri una chiamata può cambiare e quali deve conservare. \emph{Richiamo:} Una funzione chiama altre due funzioni. Quante volte deve salvare ra, e dove?
\end{chapstatus}


In questo secondo tutorial il giocatore entra nel gioco. Ogni idea usata è già
stata vista come circuito nei Capitoli da \ref{ch:boole1}
a~\ref{ch:combinational}; qui le stesse idee sono istruzioni
(Tabella~\ref{tab:v2:from}). Il file di riferimento è \fname{verm1\_part2.S}
(Appendice~\ref{app:verm1}).

\begin{table}[htb]
\centering\small
\caption{Che cosa la parte 2 prende dai capitoli sulla logica.}\label{tab:v2:from}
\begin{tabularx}{\textwidth}{@{}lLL@{}}
\toprule
\textbf{Da} & \textbf{Idea} & \textbf{Dove compare}\\
\midrule
Capitolo \ref{ch:boole1} & AND come maschera: conservare un bit, azzerare gli altri & un pulsante nella parola dei pulsanti\\
Capitolo \ref{ch:boole1} & tabelle di verità: un codice per ogni caso & le quattro direzioni in due bit\\
Capitolo \ref{ch:boole2} & XOR vale 1 dove due bit sono diversi & la regola dell'inversione a U; il test di uguaglianza\\
Capitolo \ref{ch:electronics} & un pulsante è un interruttore verso massa con un pull-up & perché «premuto» arriva come 1 solo dopo che il runtime lo ha invertito\\
Capitolo \ref{ch:combinational} & comparatore, multiplexer, ROM come tabella & la testa sul cibo; \code{dx} e \code{dy}\\
\bottomrule
\end{tabularx}
\end{table}

\section{Partite da una parte 1 che funziona}

Aprite il vostro \fname{verm1.S} della parte~1, compilatelo ed eseguitelo con i
due comandi del Listato~\ref{lst:v1:build}: la testa deve attraversare lo
schermo. Se non lo fa, copiate il file di riferimento \fname{verm1\_part1.S} e
partite da quello. Il file viene modificato in tre passi e viene eseguito dopo
il secondo e dopo il terzo. I listati segnati \textbf{modificato} sostituiscono
codice che il file già contiene; i listati segnati \textbf{nuovo} aggiungono
codice. Tenete una copia di ogni parte finita, e tenete aperto il debugger di
\prgqt: il passo~4 lo usa per guardare una decisione mentre viene presa.

% ----------------------------------------------------------------------------
\section{Passo 1: nomi per i bit}

\subsection{Costanti: pulsanti e direzioni}

Aggiungete alle costanti della parte~1 le righe di \code{RED} e quelle da
\code{BTN\_LEFT} a \code{DIR\_UP} (Listato~\ref{lst:v2:const}). C'è un colore
nuovo, il rosso, per il cibo. Ogni pulsante è un bit di una parola, e una
maschera lo seleziona. Una direzione è un numero da 0 a 3: due bit.

\asmfile[caption={Passo 1 \textbf{(modificato)}: le costanti.},label=lst:v2:const]{code/l12/c04.S}

\subsection{Un codice scelto con cura}

I codici delle quattro direzioni non sono arbitrari
(Tabella~\ref{tab:v2:dirs}). Il bit~0 dice «verticale» e il bit~1 dice
«all'indietro», così che una direzione e la sua opposta differiscono nel bit~1
e solo lì: lo XOR di un codice con il suo opposto è sempre $10_2=2$.

\begin{table}[htb]
\centering\small
\caption{Le quattro direzioni. Sullo schermo $y$ cresce verso il basso: DOWN è
$+1$, UP è $-1$.}\label{tab:v2:dirs}
\begin{tabular}{@{}lcccrrl@{}}
\toprule
\textbf{Direzione} & \textbf{Codice} & \textbf{Bit 1} & \textbf{Bit 0} & \textbf{dx} & \textbf{dy} & \textbf{Opposta}\\
\midrule
RIGHT & 0 & 0 & 0 & $+1$ & 0 & LEFT = 2\\
DOWN & 1 & 0 & 1 & 0 & $+1$ & UP = 3\\
LEFT & 2 & 1 & 0 & $-1$ & 0 & RIGHT = 0\\
UP & 3 & 1 & 1 & 0 & $-1$ & DOWN = 1\\
\bottomrule
\end{tabular}
\end{table}

\begin{prgconcept}
Una buona codifica trasforma una regola in una sola porta. È la lezione del
Capitolo~\ref{ch:boole2}, e vale più di qualsiasi trucco di programmazione.
\end{prgconcept}

\subsection{Dati: due tabelle e cinque nuove variabili}

Sostituite i dati della parte~1 con il Listato~\ref{lst:v2:data}. Le tabelle
\code{dx} e \code{dy} hanno quattro parole ciascuna, una per direzione; sono
costanti, e perciò stanno in \code{.rodata}. La variabile \code{dir} è la
direzione dell'ultimo passo, \code{want} quella chiesta dal giocatore;
\code{score}, \code{food\_x} e \code{food\_y} si usano nel passo~3.

\asmfile[caption={Passo 1 \textbf{(modificato)}: i dati.},label=lst:v2:data]{code/l12/c05.S}

% ----------------------------------------------------------------------------
\section{Passo 2: la guida}

Il piano è: leggere i pulsanti, rifiutare l'inversione a U, cercare il
movimento in una tabella.

\subsection{Dal dito al bit}

Il piedino di un pulsante legge 0 quando il pulsante è premuto: è attivo basso
(Capitolo~\ref{ch:electronics}). La funzione \code{prg32\_input\_read} inverte
e impacchetta: una parola, un bit per pulsante, con $1=\text{premuto}$.
\key{UP} e \key{RIGHT} insieme danno $\hex{04}+\hex{02}=\hex{06}$, e
$\reg{a0}\cdot\hex{04}$ vale zero se e solo se \key{UP} non è premuto. La
cartuccia non vede mai il livello elettrico: il runtime dà la stessa parola
sulla scheda, su QEMU e su \prgqt.

\subsection{\texttt{update}: dai pulsanti a un desiderio}

\asmfile[caption={Passo 2 \textbf{(modificato)}: l'inizio di \code{verm1\_update}.},label=lst:v2:wish]{code/l12/c06.S}

Il Listato~\ref{lst:v2:wish} sostituisce l'inizio di \code{verm1\_update}, che
ora chiama una funzione e quindi ha bisogno di uno stack frame. \ins{andi}
conserva un bit, e zero significa «non premuto». Ogni pulsante premuto
sovrascrive il desiderio in \reg{t2}; con due pulsanti premuti vince l'ultimo
test, così che \key{UP} batte \key{LEFT}, che batte \key{DOWN}, che batte
\key{RIGHT}. Poi viene la regola: lo XOR del desiderio e della direzione
corrente vale 2 esattamente quando sono opposti, e in quel caso il desiderio
non viene memorizzato e sopravvive quello vecchio.

Il desiderio è memorizzato in \code{want} e diventa la direzione solo al passo
successivo. In questo modo due pressioni rapide tra un passo e il successivo
non possono ripiegare il verme su sé stesso.

\begin{table}[htb]
\centering\small
\caption{La regola dell'inversione a U come tabella di verità: sei delle sedici righe.}\label{tab:v2:uturn}
\begin{tabular}{@{}llccl@{}}
\toprule
\textbf{dir} & \textbf{want} & \textbf{dir $\oplus$ want} & \textbf{$=2$?} & \textbf{Decisione}\\
\midrule
0 RIGHT & 0 RIGHT & \texttt{00} & no & prosegue\\
0 RIGHT & 1 DOWN & \texttt{01} & no & gira\\
0 RIGHT & 2 LEFT & \texttt{10} & sì & rifiuta: inversione a U\\
0 RIGHT & 3 UP & \texttt{11} & no & gira\\
1 DOWN & 3 UP & \texttt{10} & sì & rifiuta: inversione a U\\
3 UP & 2 LEFT & \texttt{01} & no & gira\\
\bottomrule
\end{tabular}
\end{table}

La regola è uno XOR e un confronto con la costante 2
(Tabella~\ref{tab:v2:uturn}). Completate le sedici righe per esercizio e
contate i rifiuti: sono quattro, uno per direzione. Senza la codifica accurata
la regola richiederebbe quattro confronti.

\subsection{\texttt{update}: il timer chiama \texttt{step}}

\asmfile[caption={Passo 2 \textbf{(modificato)}: il resto di \code{verm1\_update}.},label=lst:v2:timer]{code/l12/c07.S}

Il timer in virgola fissa della parte~1 è invariato
(Listato~\ref{lst:v2:timer}). Le righe che spostavano \code{head\_x} sono
sparite: muoversi è ora una funzione a sé, \code{step}. Le etichette locali
sono 6 e 7 perché i numeri da 1 a 5 sono usati sopra.

\subsection{\texttt{step}: cercare il movimento}

\asmfile[caption={Passo 2 \textbf{(nuovo)}: \code{step}, prima metà.},label=lst:v2:step1]{code/l12/c08.S}

Il desiderio diventa la direzione, una volta per passo
(Listato~\ref{lst:v2:step1}). Poi la direzione seleziona il movimento: una
parola è di 4 byte, quindi lo scostamento è il codice fatto scorrere a sinistra
di 2, e l'indirizzo della tabella più lo scostamento è l'indirizzo della parola
da caricare. Non c'è alcun salto: il codice seleziona, come in un multiplexer.
Una tabella in memoria di sola lettura indirizzata da un codice è una ROM
(Capitolo~\ref{ch:storage}). L'alternativa sono quattro salti; la tabella è più
corta, e cambiare le regole significa cambiare dati, non codice.

\subsection{\texttt{step}: muoversi e rientrare, nei due versi}

\asmfile[caption={Passo 2 \textbf{(nuovo)}: \code{step}, seconda metà.},label=lst:v2:step2]{code/l12/c09.S}

Nella parte~1 esisteva solo il bordo destro. Ora \code{dx} può essere $-1$ e la
colonna può scendere sotto 0: una colonna negativa è un numero con segno, e
\ins{bgez} e \ins{blt} sono i confronti con segno del
Capitolo~\ref{ch:integers} (Listato~\ref{lst:v2:step2}). Lo stesso si fa per la
riga, con \code{ROWS}. I registri \reg{t4} e \reg{t6} conservano la nuova cella
per il passo~3.

Per il momento la funzione si chiude con le tre righe del
Listato~\ref{lst:v2:end}, che ripristinano \reg{ra} e ritornano. Il passo~3
metterà prima di esse il test per il cibo.

\asmfile[caption={Passo 2: una fine provvisoria per \code{step}.},label=lst:v2:end]{code/l12/c10.S}

\begin{prgcheck}[Eseguite: state guidando]
\begin{enumerate}[nosep]
\item Compilate ed eseguite: la testa parte ancora verso destra, dalla colonna
  5 della riga 12.
\item Tenete premuto \key{DOWN}: la testa gira al suo passo successivo e
  scende, rientrando dalla riga 23 alla riga 0.
\item Mentre va a destra, tenete premuto \key{LEFT}: non succede niente.
  L'inversione a U è rifiutata.
\item Tenete premuto \key{DOWN} e poi \key{LEFT}: due svolte, e ora la testa va
  a sinistra.
\end{enumerate}
\textbf{Risultato atteso.} Tenere premuto \key{DOWN} dall'inizio per 64 frame
porta la testa in colonna 5, riga 22: dieci passi, tutti verso il basso. Tenere
premuto \key{LEFT} per 64 frame la porta in colonna 15, riga 12, come se non
fosse stato premuto niente. (Entrambi i casi sono stati verificati
sull'esecutore senza interfaccia di \prgqt.)
\end{prgcheck}

% ----------------------------------------------------------------------------
\section{Passo 3: il cibo, un punteggio, un suono}

\subsection{\texttt{init}: una partenza completa}

\asmfile[caption={Passo 3 \textbf{(modificato)}: \code{verm1\_init}.},label=lst:v2:init]{code/l12/c12.S}

Ogni variabile riceve ora il suo valore di partenza
(Listato~\ref{lst:v2:init}). \code{init} chiama una funzione, per collocare il
primo cibo, e quindi ha bisogno di uno stack frame.

\subsection{\texttt{place\_food}: una cella casuale}

\asmfile[caption={Passo 3 \textbf{(nuovo)}: \code{place\_food}.},label=lst:v2:food]{code/l12/c13.S}

La funzione \code{prg32\_random\_number(min, max)} restituisce un numero
compreso tra i suoi due argomenti, estremi inclusi; gli argomenti vanno in
\reg{a0} e \reg{a1} e il risultato torna in \reg{a0}
(Listato~\ref{lst:v2:food}). I numeri casuali vengono dal runtime, e sulla
scheda dal generatore hardware dell'ESP32-C6: due esecuzioni del gioco sono
diverse. Il cibo può cadere sulla testa; la parte~3 lo corregge.

\subsection{\texttt{step}: la testa è sul cibo?}

\asmfile[caption={Passo 3 \textbf{(modificato)}: la fine di \code{step}.},label=lst:v2:eat]{code/l12/c14.S}

Sostituite la fine provvisoria con il Listato~\ref{lst:v2:eat}. Un comparatore
di uguaglianza è una fila di porte XOR seguita da un OR
(Capitolo~\ref{ch:combinational}); qui sono due \ins{xor} e una \ins{or}: ogni
\ins{xor} vale 0 quando due numeri sono uguali, e la \ins{or} dei due risultati
vale 0 solo se entrambe le coppie sono uguali. Mangiare significa un punto in
più, nuovo cibo e un suono. I registri \reg{t4} e \reg{t6} contengono ancora la
nuova colonna e la nuova riga; dopo le chiamate non servono più, ed è per
questo che tenerli in registri temporanei è sicuro.

\subsection{\texttt{eat\_sound}: una nota}

\asmfile[caption={Passo 3 \textbf{(nuovo)}: \code{eat\_sound}.},label=lst:v2:sound]{code/l12/c15.S}

La firma è
\begin{center}\code{prg32\_audio\_note(canale, strumento, nota, volume, millisecondi)}\end{center}
con cinque argomenti, da \reg{a0} ad \reg{a4}, poi la chiamata. La nota 72 è il
Do5 nella numerazione MIDI. Su \prgqt la nota è suonata dal computer, sulla
scheda dall'amplificatore e dal suo altoparlante. Il suono viene per ultimo: il
gioco funzionava anche senza.

\subsection{\texttt{draw}: il punteggio come tacche}

\asmfile[caption={Passo 3 \textbf{(modificato)}: l'inizio di \code{verm1\_draw}.},label=lst:v2:marks]{code/l12/c16.S}

Il Listato~\ref{lst:v2:marks} disegna un piccolo rettangolo per punto, a 4
pixel di distanza, in un ciclo che termina quando il contatore raggiunge il
punteggio. Il contatore è tenuto in \reg{s0}, perché deve sopravvivere alle
chiamate: un registro \reg{t} può essere distrutto da qualsiasi chiamata,
mentre un registro \reg{s} è preservato, a condizione che la funzione che lo
usa lo salvi e lo ripristini. Quindi \reg{s0} è salvato nel frame, con
\reg{ra}. La convenzione, incontrata nel Capitolo~\ref{ch:firstcontact}, è
studiata nel Capitolo~\ref{ch:processor}. Le cifre arrivano nella parte~4,
quando sarà disponibile la divisione.

\subsection{\texttt{draw}: il cibo e la testa}

\asmfile[caption={Passo 3 \textbf{(modificato)}: il resto di \code{verm1\_draw}.},label=lst:v2:draw]{code/l12/c17.S}

Le stesse tre righe compaiono due volte, per colonna, riga e colore
(Listato~\ref{lst:v2:draw}); la conversione in pixel è lasciata a una funzione.
La testa è disegnata per ultima, così che copra il cibo.

\asmfile[caption={Passo 3 \textbf{(nuovo)}: \code{draw\_cell}.},label=lst:v2:cell]{code/l12/c18.S}

\code{draw\_cell} è l'aritmetica della parte~1, scritta una volta sola
(Listato~\ref{lst:v2:cell}). È una funzione con un contratto:
$\reg{a0}=\text{colonna}$, $\reg{a1}=\text{riga}$, $\reg{a2}=\text{colore}$.
Il colore arriva in \reg{a2} e deve uscire in \reg{a4}: viene spostato per
primo, perché \reg{a2} sta per essere sovrascritto. Scrivere il contratto in un
commento fa parte del lavoro.

\begin{prgcheck}[Eseguite: Verm1, parte 2]
Una cella rossa è da qualche parte sullo schermo. Guidateci sopra la testa:
suona una nota, compare una tacca accanto a \texttt{VERM1}, il cibo si sposta
(Figura~\ref{fig:v2:shot}).
\end{prgcheck}

\begin{figure}[htb]
\centering
\includegraphics[width=0.55\textwidth]{verm1_p2.png}
\caption{La testa, il cibo, e due tacche dopo il titolo: due punti. La
schermata è prodotta dai test, in cui un giocatore programmato guida la testa
sul cibo due volte.}
\label{fig:v2:shot}
\end{figure}

% ----------------------------------------------------------------------------
\section{Passo 4: osservare una decisione}

In questa sessione \prgqt è usato come sonda logica.

\begin{prgtry}[Cogliete la regola dell'inversione a U al lavoro]
\begin{enumerate}[nosep]
\item Eseguite la cartuccia in modalità debug e premete \key{Update}: il
  listato mostra \code{verm1\_update}.
\item Guardate la telemetria mentre tenete premuto un pulsante: la maschera di
  ingresso mostra \hex{04} per \key{UP}, \hex{01} per \key{LEFT}.
\item Trovate l'istruzione \code{xor t3, t2, t1} e fate clic sul margine: un
  breakpoint.
\item Tenete premuto \key{LEFT} mentre la testa va a destra, e premete
  \key{Resume}: il programma si ferma sulla \ins{xor}.
\item Leggete i registri: \reg{t1} vale 0, la direzione; \reg{t2} vale 2, il
  desiderio.
\item Premete \key{Step}: \reg{t3} diventa 2. Fate altri due passi: il salto
  viene preso e la \ins{sw} viene saltata.
\end{enumerate}
Abbassate la velocità del debugger a 0.25x se vi serve tempo per premere.
\end{prgtry}

La Tabella~\ref{tab:v2:regs} elenca ciò che si legge nei registri, sempre con
la testa che va a destra. La sua ultima riga mostra la priorità dei test:
\key{UP} è esaminato per ultimo e vince. \reg{a0} viene sovrascritto più avanti
nel frame, e va letto alla prima istruzione dopo la chiamata.

\begin{table}[htb]
\centering\small
\caption{Che cosa mostrano i registri al breakpoint.}\label{tab:v2:regs}
\begin{tabular}{@{}lccccl@{}}
\toprule
\textbf{Tenete premuto} & \textbf{\reg{a0}} & \textbf{\reg{t1} (dir)} & \textbf{\reg{t2} (want)} & \textbf{\reg{t3} = \reg{t2}$\oplus$\reg{t1}} & \textbf{Memorizzato?}\\
\midrule
niente & \hex{00} & 0 & 0 & 0 & sì: 0\\
\key{DOWN} & \hex{08} & 0 & 1 & 1 & sì: 1\\
\key{LEFT} & \hex{01} & 0 & 2 & 2 & no: inversione a U\\
\key{UP} & \hex{04} & 0 & 3 & 3 & sì: 3\\
\key{UP} + \key{RIGHT} & \hex{06} & 0 & 3 & 3 & sì: 3\\
\bottomrule
\end{tabular}
\end{table}

\subsection{Leggere le variabili in memoria}

Dopo \code{la t0, dir} il registro \reg{t0} contiene l'indirizzo di
\code{dir}: scrivetelo nel campo \emph{Memory}. Le variabili si susseguono, 4
byte ciascuna, nell'ordine del file: dall'indirizzo di \code{dir}, a $+0$ c'è
\code{dir}, a $+4$ \code{want}, a $+8$ \code{step\_q8}, a $+12$ \code{score}, a
$+16$ \code{food\_x} e a $+20$ \code{food\_y}; \code{head\_x} e \code{head\_y}
sono le due parole prima di \code{dir}. Mangiate una volta, mettete in pausa, e
leggete \code{score}: \texttt{01 00 00 00}, il byte meno significativo per
primo. Poi leggete \code{food\_x} e \code{food\_y} e confrontateli con lo
schermo: sono la colonna e la riga della cella rossa.

% ----------------------------------------------------------------------------
\section{Verificate di aver capito}

\begin{quiz}{Il giocatore tiene premuti insieme \key{DOWN} e \key{RIGHT}: \code{prg32\_input\_read} restituisce}
\item \hex{02}
\item \hex{08}
\item \hex{0a}
\item \hex{10}
\end{quiz}

\begin{quiz}{Il verme va verso UP (3). Quale desiderio viene rifiutato?}
\item RIGHT (0)
\item DOWN (1)
\item LEFT (2)
\item UP (3)
\end{quiz}

\begin{quiz}{Perché \code{slli t1, t1, 2} prima di sommare all'indirizzo di \code{dx}?}
\item per rendere positivo il codice
\item ogni elemento della tabella è di 4 byte
\item per azzerare il bit 1
\item le tabelle cominciano all'indirizzo 4
\end{quiz}

\begin{quiz}{\code{xor t1, t1, t4} dà zero esattamente quando}
\item \reg{t1} vale zero
\item \reg{t4} vale zero
\item \reg{t1} e \reg{t4} sono uguali
\item \reg{t1} e \reg{t4} sono opposti
\end{quiz}

\section{Fatelo vostro}

\begin{exercise}[su \prg]
\begin{enumerate}[nosep]
\item Fate del pulsante \key{A} un freno: mentre è tenuto premuto, il timer non avanza.
\item Date due punti per un pezzo di cibo mangiato nella riga più in alto.
\item Disegnate la testa in un colore diverso per ogni direzione, usando una
  terza tabella di quattro colori.
\item Cambiate la codifica in modo che RIGHT = 0, LEFT = 1, DOWN = 2, UP = 3:
  che cosa diventa la regola dell'inversione a U?
\end{enumerate}
\expected{Ogni modifica è di poche righe. Tenete una copia della parte~2
funzionante prima di sperimentare: la parte~3 parte da lì. (Suggerimenti: per
il freno, \ins{andi} con \code{BTN\_A}, \hex{10}, all'inizio del timer e un
salto all'uscita. La tabella dei colori è fatta di quattro \code{.word} in
\code{.rodata}, indicizzata come \code{dx}. Con la nuova codifica gli opposti
differiscono nel bit~0, quindi il test diventa $\oplus = 1$.)}
\end{exercise}

\section{Dove siamo, e che cosa viene dopo}

Il giocatore è nel gioco: un desiderio, una regola, un movimento, una
ricompensa. Bit e maschere, XOR come comparatore, una tabella come selettore:
la logica è stata trasformata in istruzioni. Ma il verme non ha un corpo,
niente può andare storto, e il gioco non finisce mai. I Capitoli da
\ref{ch:sequential} a~\ref{ch:storage} danno gli strumenti: macchine a stati,
registri, la memoria come array, e ciò che sopravvive quando si toglie
l'alimentazione. La parte~3 aggiunge un corpo che cresce, le collisioni, una
schermata del titolo e una fine partita.
